\documentclass[11pt, a4paper]{article}
\usepackage{geometry}
\usepackage[parfill]{parskip}
\usepackage{graphicx}
\usepackage{amsmath}
\usepackage{enumerate}

\title{Huiswerk 3 \\ IB0402 Logica, verzamelingen en relaties}
\author{Harman Singh}
\date{Month Day, Year}

\begin{document}
\maketitle

\pagebreak
\section*{Opgave 7}

\textbf{A}

Aangezien een machtsverzameling $2^n$ elementen bevat, en B 3 elementen bevat, zal $\mathcal{P(B)}$ dus $2^3 = 8$ elementen bevatten.

Namelijk:

\begin{itemize}
	\item $\emptyset$ \hfill(lege verzameling)
	\item $\{a\}$ \hfill(het 1e element)
	\item $\{\{b,a\}\}$ \hfill(het 2e element)
	\item $\{b\}$ \hfill(het 3e element)
	\item $\{a,b\}$ \hfill(2e en 3e element)
	\item $\{a,\{b,a\}\}$ \hfill(1e en 2e element)
	\item $\{\{b,a\},b\}$ \hfill(2e en 3e element)
	\item $\{a, \{b,a\},b\}$ \hfill(de hele verzameling)
\end{itemize}

\[
\boxed{\mathcal{P(B)} = \{\emptyset, \{a\}, \{\{b,a\}\}, \{b\}, \{a,b\}, \{a,\{b,a\}\}, \{\{b,a\},b\}, \{a, \{b,a\},b\}\}}
\]


\textbf{B}

Nu dat we $\mathcal{P(B)}$ kennen, moeten we voor $A \cap \mathcal{P(B)}$ enkel de elementen overhouden die in beide verzamelingen voorkomen.

En dat is enkel: $\{a,b\}$

\[\boxed{A \cap \mathcal{P(B)} = \{a,b\}}\]

\textbf{C}

Eerst zoek in naar alle verzamelingen in B, en dat is enkel $\{b,a\}$. Dan moet ik kijken of alle elementen van $\{b,a\}$ ook in A zitten. Dat is niet het geval, dus het antwoord is een lege verzameling: $\emptyset$.

\[\boxed{B \cap \mathcal{P}(A) = \emptyset} \]



\pagebreak
\section*{Opgave 8}
\ldots

\section*{Opgave 9}
\ldots

\section*{Opgave 10}
\ldots


\ldots
\end{document}